\documentclass[11pt]{article}
\usepackage{fontspec}
\setmainfont{Times New Roman}
\setsansfont{Arial}
\setmonofont{Consolas}
\usepackage{amsmath,amssymb,mathtools}
\usepackage{geometry}
\usepackage{microtype}
\newcommand{\mo}{\mathfrak{o}}
\newcommand{\mO}{\mathfrak{O}}
\newcommand{\mT}{\mathfrak{T}}
\geometry{a4paper,margin=2.35cm}
\setlength{\parindent}{1.25em}
\setlength{\parskip}{0pt}
\makeatletter
\renewcommand\footnotesize{\@setfontsize\footnotesize{7.7}{8.7}\selectfont}
\makeatother
\emergencystretch=100em
\allowdisplaybreaks
\sloppy
\hbadness=10000
\vbadness=10000
\hfuzz=2pt
\vfuzz=10pt
\raggedbottom

\begin{document}

% Унит: Папер31 Сецтион07 но. 1 опенинг в001. Дирецт Интерславиц Латин транслатион фром the Герман финал-аудитед соурце слице, соурце линес 385--389 / сегментс P31-С0082--P31-С0084, експандед агаинст принтед сцан/ОЦР витнесс пагес 98--99.
\setcounter{footnote}{28}

\subsection*{§7. Теорем о дискриминанту за порјадкы алгебраичного числового польа или функционалного польа}

\noindent\textbf{1.}\quad Разсматриване числове и функционалне польа -- при чем при алгебраичных функцијах од веч неж једнеј неозначене и при целых алгебраичных функцијах иде се о разширенье функционалној области -- все подлагајут се следујучему типу польа с фиксованым појмом целых елементов. Нехај $\mo$ буде колцо главных идеалов без делительев нулы и с јединицеју (колцо целых рационалных чисел, област полиномов од једнеј неозначене с коефициентами из польа, функционална област), а $\Omega$ јего полье квоциентов. Нехај $K$ буде конечно разширенье првого рода польа $\Omega$, а $\mO$ колцо всих елементов из $K$, целых взходно к $\mo$. Всако подколцо од $\mO$, кторо обејмаје $\mo$, называ се порјадок; само $\mO$ означаје се како главны порјадок.

Всакы $\mo$-модул в $\mO$, особливо все идеалы порјадка и сам порјадок, имаје линеарно независну модулну базу взходно к $\mo$. Ако $K$ имаје степен $n$ над $\Omega$, тогда доста је ограничити се на порјадкы ранга $n$ взходно к $\mo$, бо всакы порјадок нижшего ранга чрез творјенье квоциентов пораждаје меджуполье меджу $K$ и $\Omega$ именно тој степени, и за јего все предпосылкы оставајут. Зато всакы порјадок $\mT$ јест колцо, за кторо предпосылкы §4 и §5 сут исполньене, и именно колцо без делительев нулы; тым дискриминант $D_\mT$ порјадка јест једнозначно определьен до фактора, кторы јест квадрат јединице из $\mo$. Та дискриминант јест такоже, до фактора-јединице из $\Omega$, идентичны с дискриминантом польа $K$, сматрјаным како $\Omega$-модул, и зато различны од нулы, бо $K$ был предположены како разширенье првого рода (§6, 3). Дискриминант такоже допушча представјенье, дано в §6, 4, чрез квадрат детерминанта конјугованых базных елементов.

\end{document}

